% physics_topological_strings.tex
% Physics note: Topological string theory on CY3 and the root datum
% Raeez Lorgat, April 2026
%
% This note develops the dictionary between the A/B topological string
% on a Calabi--Yau threefold X and the quantum vertex chiral group G(X)
% of Volume III.  The four pillars: (1) A-model and GW invariants as
% genus-g shadow amplitudes; (2) B-model holomorphic anomaly = MC equation;
% (3) topological vertex as local root datum; (4) connection to the
% Vol I shadow obstruction tower via F_g = kappa * lambda_g.

\documentclass[11pt]{amsart}
\usepackage[T1]{fontenc}
\usepackage[utf8]{inputenc}
\usepackage{amsmath,amssymb,amsthm}
\usepackage{mathtools}
\usepackage{enumitem}
\usepackage[margin=1.25in]{geometry}
\usepackage{hyperref}

\theoremstyle{plain}
\newtheorem{theorem}{Theorem}[section]
\newtheorem{proposition}[theorem]{Proposition}
\newtheorem{lemma}[theorem]{Lemma}
\newtheorem{corollary}[theorem]{Corollary}
\newtheorem{conjecture}[theorem]{Conjecture}
\theoremstyle{definition}
\newtheorem{definition}[theorem]{Definition}
\newtheorem{example}[theorem]{Example}
\newtheorem{construction}[theorem]{Construction}
\theoremstyle{remark}
\newtheorem{remark}[theorem]{Remark}
\newtheorem{warning}[theorem]{Warning}

% Macros (matching Volume III main.tex conventions)
\newcommand{\CY}{\mathrm{CY}}
\newcommand{\HH}{\mathrm{HH}}
\newcommand{\HC}{\mathrm{HC}}
\newcommand{\Tr}{\mathrm{Tr}}
\newcommand{\Ran}{\mathrm{Ran}}
\newcommand{\Fact}{\mathrm{Fact}}
\newcommand{\Rep}{\mathrm{Rep}}
\newcommand{\Fuk}{\mathrm{Fuk}}
\newcommand{\Coh}{\mathrm{Coh}}
\newcommand{\Mod}{\mathrm{Mod}}
\newcommand{\Sym}{\mathrm{Sym}}
\newcommand{\id}{\mathrm{id}}
\newcommand{\Ainf}{A_\infty}
\newcommand{\cC}{\mathcal{C}}
\newcommand{\cM}{\mathcal{M}}
\newcommand{\cO}{\mathcal{O}}
\newcommand{\cF}{\mathcal{F}}
\newcommand{\cZ}{\mathcal{Z}}
\newcommand{\frakg}{\mathfrak{g}}
\renewcommand{\hbar}{\hslash}
\newcommand{\Mbar}{\overline{\mathcal{M}}}
\newcommand{\DT}{\mathrm{DT}}
\newcommand{\GW}{\mathrm{GW}}
\newcommand{\GV}{\mathrm{GV}}

\title{Topological String Theory on Calabi--Yau Threefolds\\
  and the Root Datum of the Quantum Vertex Chiral Group}
\author{Raeez Lorgat}
\date{April 2026}

\begin{document}
\maketitle

\begin{abstract}
We develop the dictionary between topological string theory on a Calabi--Yau
threefold $X$ and the quantum vertex chiral group $G(X)$ of Volume~III.
The A-model free energy $F = \sum_g F_g\, g_s^{2g-2}$, with $F_g$
encoding Gromov--Witten invariants $N_{g,\beta}$, is identified with the genus
expansion of the shadow obstruction tower $\Theta_{A_X}$.  The B-model holomorphic
anomaly equation of BCOV is reinterpreted as the Maurer--Cartan equation
$D\Theta + \frac{1}{2}[\Theta, \Theta] = 0$ projected to genus~$g$.
The topological vertex (AKMV) for toric CY3 is shown to factorize as a
tensor product of local root data, one per vertex of the toric diagram.
The leading genus-$g$ term $F_g = \kappa \cdot \lambda_g$ is the
uniform-weight lane of the Vol~I shadow obstruction tower; higher terms are shadow
corrections from imaginary roots.
\end{abstract}

\tableofcontents

%% ================================================================
\section{Introduction and overview}
\label{sec:intro}

Let $X$ be a compact (or appropriately completed) Calabi--Yau threefold.
The topological string on $X$ comes in two versions --- the A-model and the
B-model --- related by mirror symmetry $X \leftrightarrow X^\vee$.  The
central object on both sides is the free energy
\begin{equation}\label{eq:free-energy}
  F(g_s; t) = \sum_{g=0}^{\infty} F_g(t)\, g_s^{2g-2},
\end{equation}
where $g_s$ is the string coupling and $t$ denotes moduli (K\"ahler on the
A-side, complex structure on the B-side).  The genus-$g$ amplitude $F_g$ is
the string-theoretic analogue of the genus-$g$ obstruction
$\mathrm{obs}_g(A_X)$ of the quantum vertex chiral group $G(X)$ constructed
in Volume~III.

The purpose of this note is to make the following identifications precise:

\begin{center}
\renewcommand{\arraystretch}{1.5}
\begin{tabular}{ll}
  \hline
  \textbf{Topological string on $X$} & \textbf{Quantum vertex chiral group $G(X)$} \\
  \hline
  Free energy $F$ & Shadow obstruction tower $\Theta_{A_X}$ \\
  Genus-$g$ amplitude $F_g$ & Genus-$g$ obstruction $\mathrm{obs}_g(A_X)$ \\
  GW invariants $N_{g,\beta}$ & Shadow amplitudes of $G(X)$ at genus $g$ \\
  BCOV holomorphic anomaly & MC equation $D\Theta + \tfrac{1}{2}[\Theta,\Theta] = 0$ \\
  Topological vertex $C_{\lambda\mu\nu}(q)$ & Local root datum at a toric vertex \\
  DT partition function $Z^{\DT}(X)$ & Bar-complex Euler product of $G(X)$ \\
  Leading $F_g$ on uniform lane & $\kappa(G(X)) \cdot \lambda_g$ \\
  \hline
\end{tabular}
\end{center}

\noindent
The identification $g_s = \hbar$ (string coupling $=$ the quantization
parameter of the chiral algebra) and $Q^\beta = e^{-t \cdot \beta}$
(exponentiated K\"ahler class $=$ root-lattice fugacity) will be implicit
throughout.


%% ================================================================
\section{The A-model topological string on CY3}
\label{sec:a-model}

\subsection{The A-model partition function}
\label{sub:a-model-partition}

The A-model topological string on a CY3 $X$ localizes onto holomorphic maps
$f \colon \Sigma_g \to X$ from a genus-$g$ worldsheet to $X$.  The genus-$g$
free energy is a generating function for Gromov--Witten invariants:
\begin{equation}\label{eq:fg-a}
  F_g(t) = \sum_{\beta \in H_2(X,\mathbb{Z})} N_{g,\beta}\, Q^\beta,
\end{equation}
where $Q^\beta = e^{-\int_\beta \omega}$ for $\omega$ the complexified
K\"ahler form, and $N_{g,\beta}$ is the genus-$g$ degree-$\beta$
Gromov--Witten invariant:
\[
  N_{g,\beta} = \int_{[\Mbar_{g,0}(X,\beta)]^{\mathrm{vir}}} 1.
\]
The virtual fundamental class $[\Mbar_{g,0}(X,\beta)]^{\mathrm{vir}}$ has
virtual dimension zero for CY3 (since $c_1(X) = 0$), so $N_{g,\beta}$ is
a number --- a virtual count of genus-$g$ holomorphic curves in class
$\beta$.

The total free energy
\[
  F = \sum_{g=0}^{\infty} g_s^{2g-2} \sum_{\beta} N_{g,\beta}\, Q^\beta
\]
is the A-model topological string partition function (at the level of the
free energy; the partition function proper is $Z = e^F$).

\subsection{Gopakumar--Vafa reformulation}
\label{sub:gv}

The Gopakumar--Vafa conjecture (now theorem, Ionel--Parker) repackages the GW
invariants into integer BPS invariants $n^g_\beta \in \mathbb{Z}$:
\begin{equation}\label{eq:gv}
  F = \sum_{g=0}^{\infty} \sum_{\beta \neq 0} \sum_{k=1}^{\infty}
    \frac{n^g_\beta}{k}
    \left(2\sin\frac{k g_s}{2}\right)^{2g-2} Q^{k\beta}.
\end{equation}
The integers $n^g_\beta$ count BPS states of charge $\beta$ and spin $g$ in
the M-theory lift.  In the language of the quantum vertex chiral group:
\begin{itemize}
  \item $n^0_\beta$ are the genus-0 BPS counts $=$ DT invariants
    $=$ root multiplicities $\mathrm{mult}(\beta)$ of $G(X)$.
  \item $n^g_\beta$ for $g > 0$ are the genus-$g$ refinements of the root
    multiplicity, living in the $g$-th shadow correction.
\end{itemize}

\begin{remark}[GV invariants as root data]
\label{rem:gv-root}
The generating function $\sum_\beta n^0_\beta\, Q^\beta$ is the
multiplicity-generating function of the imaginary roots of $G(X)$.  For the
$K3 \times E$ family, these are the Fourier coefficients of $\phi_{0,1}$
(the K3 elliptic genus).  For a toric CY3, they are the DT invariants of the
quiver with potential $(Q_X, W_X)$.  The higher-genus invariants $n^g_\beta$
refine the root datum by incorporating information from the Hodge bundle
$\mathbb{E}_g \to \Mbar_{g,0}$; in the shadow obstruction tower language, they appear
as genus-$g$ shadow corrections to the tree-level root multiplicities.
\end{remark}

\subsection{GW invariants as shadow amplitudes}
\label{sub:gw-shadow}

The quantum vertex chiral group $G(X) = \Phi(\cC_X)$ (where $\cC_X$ is the
appropriate CY category of $X$ --- the Fukaya category on the A-side, or
$D^b(\Coh(X^\vee))$ on the B-side via HMS) carries a shadow obstruction tower
$\Theta_{A_X}$ (Volume~I).  The genus expansion of $\Theta_{A_X}$ is:
\begin{equation}\label{eq:theta-genus}
  \Theta_{A_X} = \sum_{g=0}^{\infty} \hbar^{2g-2}\, \Theta^{(g)}_{A_X},
\end{equation}
where $\Theta^{(g)}_{A_X} \in
C^\bullet(\Mbar_{g,n},\, B(A_X)^{\otimes n})$ are the genus-$g$ shadow
amplitudes --- sections of the $n$-fold tensor product of the bar complex
over the moduli space of $n$-pointed genus-$g$ curves.

\begin{conjecture}[GW/shadow identification]
\label{conj:gw-shadow}
For a CY3 $X$ and its quantum vertex chiral group $G(X)$:
\begin{enumerate}[label=(\roman*)]
  \item The genus-$g$, $n$-point shadow amplitude
    $\Theta^{(g)}_{A_X}|_{\beta}$ evaluated on the class $\beta$ recovers the
    GW descendant invariant:
    \[
      \langle \Theta^{(g)}_{A_X} \rangle_\beta
      = \int_{[\Mbar_{g,n}(X,\beta)]^{\mathrm{vir}}} \psi_1^{a_1}
        \cdots \psi_n^{a_n}.
    \]
  \item The genus-$g$ free energy $F_g$ is the $0$-point specialization:
    \[
      F_g = \langle \Theta^{(g)}_{A_X} \rangle
      = \mathrm{obs}_g(A_X).
    \]
  \item The full topological string free energy is the shadow obstruction tower:
    \[
      F = \Theta_{A_X}
    \]
    under the identification $g_s = \hbar$, $Q^\beta = e^{-t \cdot \beta}$.
\end{enumerate}
\end{conjecture}

\begin{remark}
The identification $F = \Theta_{A_X}$ is of course schematic: $F$ is a
formal power series in $g_s$ and $Q^\beta$, while $\Theta_{A_X}$ is a
Maurer--Cartan element in the dgLa governing deformations of the chiral bar
complex.  The precise statement is that the formal expansion of $\Theta_{A_X}$
in the genus grading and the root-lattice grading reproduces $F$.
\end{remark}


%% ================================================================
\section{The B-model and the holomorphic anomaly equation}
\label{sec:b-model}

\subsection{Special geometry of the CY moduli space}
\label{sub:special-geometry}

Let $X^\vee$ be the mirror CY3 and $\cM = \cM_{\mathrm{cpx}}(X^\vee)$ its
complex structure moduli space, of complex dimension $h = h^{2,1}(X^\vee)$.
The B-model free energy $F_g$ is a function (more precisely, a section of a
line bundle) on $\cM$.  The moduli space $\cM$ carries the special K\"ahler
structure:
\begin{itemize}
  \item Special coordinates $t^i$, $i = 1, \ldots, h$;
  \item The prepotential $\cF(t) = F_0(t)$ (the genus-0 free energy);
  \item The period matrix $\tau_{ij} = \partial_i \partial_j \cF$;
  \item The Weil--Petersson metric $G_{i\bar{j}} = -\partial_i \partial_{\bar{j}} K$
    where $K$ is the K\"ahler potential;
  \item The Yukawa coupling $C_{ijk} = \partial_i \partial_j \partial_k \cF$
    (the genus-0 three-point function);
  \item The propagator $S^{ij}$ satisfying $\partial_{\bar{k}} S^{ij} = \bar{C}^{ij}_{\bar{k}} = e^{2K} G^{i\bar{i}'} G^{j\bar{j}'} \bar{C}_{\bar{k}\bar{i}'\bar{j}'}$.
\end{itemize}
The propagator $S^{ij}$ is the genus-1 annulus amplitude: it mediates the
degeneration of a genus-$g$ surface into lower genera by pinching a cycle.

\subsection{The BCOV holomorphic anomaly equation}
\label{sub:bcov}

The central result of Bershadsky--Cecotti--Ooguri--Vafa (BCOV) is that the
genus-$g$ free energy, while not holomorphic, satisfies a controlled
recursion in $g$:

\begin{equation}\label{eq:bcov}
  \partial_{\bar{i}} F_g = \frac{1}{2} \bar{C}^{jk}_{\bar{i}}
    \left( D_j D_k F_{g-1} + \sum_{r=1}^{g-1} D_j F_r\, D_k F_{g-r} \right),
  \qquad g \geq 2.
\end{equation}

Here $D_i = \nabla_i + (2-2g) \partial_i K$ is the covariant derivative
incorporating the K\"ahler connection, and $\bar{C}^{jk}_{\bar{i}} =
e^{2K} G^{j\bar{j}'} G^{k\bar{k}'} \bar{C}_{\bar{i}\bar{j}'\bar{k}'}$
is the anti-holomorphic Yukawa coupling raised by the metric.  This equation
says: the failure of $F_g$ to be holomorphic is controlled by the
propagator and the lower-genus data.

\begin{remark}[The genus-1 anomaly]
At genus $g = 1$, the anomaly takes a different form:
\[
  \partial_{\bar{i}} \partial_j F_1 = \frac{1}{2} C_{jkl}\, \bar{C}^{kl}_{\bar{i}}
    - \left(\frac{\chi(X)}{24} - 1\right) G_{j\bar{i}},
\]
where $\chi(X)$ is the Euler characteristic.  The constant $\chi(X)/24$
is precisely the modular characteristic $\kappa(A_X)$ from Vol~I: it
controls the curvature of the Hodge line bundle $\lambda_1$ on $\Mbar_{1,0}$.
\end{remark}

\subsection{BCOV as Maurer--Cartan}
\label{sub:bcov-mc}

The key insight, which connects the B-model to the framework of Volume~I, is
that the BCOV holomorphic anomaly equation \eqref{eq:bcov} IS the
Maurer--Cartan equation
\begin{equation}\label{eq:mc}
  D\Theta + \frac{1}{2}[\Theta, \Theta] = 0
\end{equation}
projected to genus $g$.

\begin{construction}[BCOV as MC]
\label{constr:bcov-mc}
The dictionary is as follows.

\begin{enumerate}[label=(\roman*)]
  \item \textbf{The dgLa.}  The relevant dgLa is the BCOV complex
    \[
      \mathcal{L}_X = \mathrm{PV}^{\bullet,\bullet}(X^\vee) =
      \bigoplus_{p,q} \Omega^{0,q}(X^\vee, \bigwedge^p TX^\vee),
    \]
    the space of polyvector-valued $(0,q)$-forms on $X^\vee$, with
    differential $\bar{\partial}$ and bracket given by the Schouten--Nijenhuis
    bracket wedged with the Dolbeault differential.  This is the
    Kodaira--Spencer dgLa governing deformations of $X^\vee$.

  \item \textbf{The MC element.}  The total free energy
    \[
      \Theta = \sum_{g=0}^{\infty} \hbar^{2g-2}\, \Theta^{(g)},
      \qquad \Theta^{(g)} \sim F_g,
    \]
    assembled from all genera, is a MC element of $\mathcal{L}_X[[\hbar]]$.

  \item \textbf{The genus-$g$ projection.}  The MC equation
    $D\Theta + \frac{1}{2}[\Theta,\Theta] = 0$, expanded in powers of
    $\hbar^{2g-2}$, gives at genus $g$:
    \begin{equation}\label{eq:mc-genus-g}
      D\Theta^{(g)} + \frac{1}{2} \sum_{r=1}^{g-1}
      [\Theta^{(r)}, \Theta^{(g-r)}] = 0.
    \end{equation}
    This is precisely the BCOV equation \eqref{eq:bcov}: the differential
    $D$ encodes $\bar{\partial}$ plus the K\"ahler connection, while
    the bracket $[\Theta^{(r)}, \Theta^{(g-r)}]$ corresponds to the
    bilinear term $D_j F_r\, D_k F_{g-r}$ contracted with the propagator
    $\bar{C}^{jk}_{\bar{i}}$.

  \item \textbf{The propagator as bracket.}  The propagator $S^{ij}$
    of special geometry is the homotopy transfer data for the
    Hodge-to-de-Rham spectral sequence on $\mathrm{PV}^{\bullet,\bullet}$.
    Concretely:
    \[
      [\Theta^{(r)}, \Theta^{(g-r)}]_{\bar{i}}
      = \bar{C}^{jk}_{\bar{i}}\, D_j \Theta^{(r)} \cdot D_k \Theta^{(g-r)},
    \]
    so the Lie bracket in $\mathcal{L}_X$ is ``propagator-dressed.''

  \item \textbf{The non-holomorphic completion.}  The fact that $\Theta^{(g)}$
    is not holomorphic (the anomaly) is the statement that $\Theta$ is a MC
    element of the full dgLa $\mathcal{L}_X$, not of the holomorphic
    sub-dgLa $\mathrm{PV}^{\bullet,0}$.  Holomorphic ambiguity (the freedom
    in completing $F_g$ to a modular object) corresponds to the freedom in
    choosing a representative within the MC gauge-equivalence class.
\end{enumerate}
\end{construction}

\begin{proposition}[MC = BCOV, schematic]
\label{prop:mc-bcov}
The Maurer--Cartan equation $D\Theta + \frac{1}{2}[\Theta,\Theta] = 0$ in the
dgLa $\mathcal{L}_X = \mathrm{PV}^{\bullet,\bullet}(X^\vee)[[\hbar]]$, expanded
at order $\hbar^{2g-2}$, reproduces the BCOV holomorphic anomaly equation
\eqref{eq:bcov}.  The propagator $S^{ij}$ of special geometry is the
contraction kernel of the Lie bracket; the Yukawa coupling $C_{ijk}$ is the
structure constant of the Schouten--Nijenhuis bracket; and the K\"ahler
covariant derivative $D_i$ is the dgLa differential restricted to the
holomorphic tangent directions.
\end{proposition}

\begin{remark}
This is, of course, the perspective championed by Costello--Li, who construct
the BCOV theory as a quantization of the Kodaira--Spencer dgLa.  What we add
is the identification of this dgLa with the chiral bar dgLa of
Volume~I: the BCOV complex $\mathcal{L}_X$ is quasi-isomorphic to the
dgLa $C^\bullet(B(A_X))$ governing deformations of the bar complex of the
quantum vertex chiral group $G(X)$.  In this identification, the BCOV
MC element $\Theta$ IS the shadow obstruction tower $\Theta_{A_X}$.
\end{remark}

\subsection{The holomorphic anomaly and the shadow obstruction tower}
\label{sub:anomaly-shadow}

The holomorphic anomaly has a clean interpretation in the shadow obstruction tower:

\begin{center}
\renewcommand{\arraystretch}{1.5}
\begin{tabular}{ll}
  \hline
  \textbf{B-model} & \textbf{Shadow obstruction tower (Vol~I)} \\
  \hline
  Holomorphic limit ($\bar{t} \to \infty$) & Tree-level (genus 0) shadow \\
  Anti-holomorphic dependence of $F_g$ & Genus-$g$ MC obstruction \\
  Propagator $S^{ij}$ & Homotopy transfer / degeneration data \\
  Yukawa $C_{ijk}$ & Cubic shadow $C$ (arity 3) \\
  Holomorphic ambiguity & MC gauge freedom \\
  Modular completion & Full automorphic correction \\
  \hline
\end{tabular}
\end{center}

\noindent
The ``holomorphic limit'' of the B-model (sending $\bar{t} \to \infty$,
i.e.\ restricting to the large-complex-structure point) corresponds to
retaining only the tree-level shadow: the genus-0 data, which is purely
holomorphic.  Each genus $g \geq 1$ introduces anti-holomorphic dependence,
which in the shadow obstruction tower corresponds to the $g$-th obstruction class
living in a cohomology group that is not concentrated in the holomorphic
sector.


%% ================================================================
\section{The topological vertex for toric CY3}
\label{sec:top-vertex}

\subsection{Toric CY3 and their diagrams}
\label{sub:toric-cy3}

A toric CY3 $X_\Sigma$ is determined by a fan $\Sigma$ in $\mathbb{Z}^3$
with all generators coplanar (the CY condition).  The toric diagram is the
projection of the fan to a convex polygon in $\mathbb{Z}^2$.  The geometry
of $X_\Sigma$ is encoded in the trivalent planar graph dual to the toric
diagram: vertices correspond to $\mathbb{C}^3$ patches, edges correspond to
$\mathbb{C}^2$ patches (``gluing strips''), and the external legs correspond
to non-compact divisors.

\begin{example}[$\mathbb{C}^3$: the simplest vertex]
A single vertex.  $X = \mathbb{C}^3$, quiver = Jordan quiver, $G(X) \supset Y(\widehat{\mathfrak{gl}}_1) \simeq \mathcal{W}_{1+\infty}$.
\end{example}

\begin{example}[Resolved conifold]
Two vertices joined by an internal edge.  $X = \mathcal{O}(-1) \oplus \mathcal{O}(-1) \to \mathbb{P}^1$.  The partition function $Z = M(q)^2 \sum_\lambda q^{|\lambda|} Q^{|\lambda|} (s_\lambda(q^\rho))^2$ where $M(q) = \prod_{n=1}^\infty (1-q^n)^{-n}$ is the MacMahon function.
\end{example}

\subsection{The AKMV topological vertex}
\label{sub:akmv}

The fundamental result of Aganagic--Klemm--Mari\~no--Vafa (AKMV) is that the
A-model topological string partition function on any toric CY3 $X$
\emph{factorizes} over the vertices and edges of the toric diagram:

\begin{equation}\label{eq:top-vertex-Z}
  Z(X) = \sum_{\{\lambda_e\}} \prod_{\text{vertices } v}
    C_{\lambda_{e_1} \lambda_{e_2} \lambda_{e_3}}(q)
    \prod_{\text{edges } e} (-Q_e)^{|\lambda_e|}\, f_{\lambda_e}(q),
\end{equation}
where:
\begin{itemize}
  \item The sum is over assignments of partitions $\lambda_e$ to each internal
    edge $e$ of the toric diagram;
  \item $C_{\lambda\mu\nu}(q)$ is the \emph{topological vertex} --- a
    universal amplitude attached to each trivalent vertex, depending on the
    three partitions assigned to its edges;
  \item $Q_e = e^{-t_e}$ is the K\"ahler parameter of the edge $e$ (the
    exponentiated area of the corresponding $\mathbb{P}^1$);
  \item $f_{\lambda}(q) = (-1)^{|\lambda|} q^{(\kappa_\lambda)/2}$ is the
    framing factor, with $\kappa_\lambda = \sum_i \lambda_i(\lambda_i - 2i + 1)$.
\end{itemize}

The topological vertex itself has the explicit expression (in the
Schur-function basis):
\begin{equation}\label{eq:c-vertex}
  C_{\lambda\mu\nu}(q) = q^{\frac{\kappa_\mu}{2}}\, s_{\nu^t}(q^\rho)
    \sum_\eta s_{\lambda^t/\eta}(q^{\nu + \rho})\, s_{\mu/\eta}(q^{\nu^t + \rho}),
\end{equation}
where $q^\rho = (q^{-1/2}, q^{-3/2}, q^{-5/2}, \ldots)$ and
$q^{\nu + \rho} = (q^{\nu_1 - 1/2}, q^{\nu_2 - 3/2}, \ldots)$ are
specializations of symmetric functions, and $s_{\lambda/\eta}$ are skew Schur
functions.

\subsection{The topological vertex as local root datum}
\label{sub:vertex-root-datum}

The topological vertex $C_{\lambda\mu\nu}(q)$ is a ``local root datum'' of
$G(X)$.  The precise formulation:

\begin{construction}[Local root datum]
\label{constr:local-root-datum}
To each trivalent vertex $v$ of the toric diagram, associate the local
quantum vertex chiral group $G_v = G(\mathbb{C}^3)$, with:
\begin{enumerate}[label=(\roman*)]
  \item \textbf{Lattice}: $\Lambda_v = K_0(\mathrm{pt}) = \mathbb{Z}$
    (the dimension vector).
  \item \textbf{Root datum}: the affine Yangian $Y(\widehat{\mathfrak{gl}}_1)$
    root system.  The root multiplicities $\mathrm{mult}_v(n) = p(n)$
    (partitions of $n$), controlled by the MacMahon function $M(q)$.
  \item \textbf{Vertex amplitude}: the topological vertex
    $C_{\lambda\mu\nu}(q)$ is the matrix element of the intertwiner
    \[
      \Phi_v \colon \cF_{\lambda} \otimes \cF_{\mu} \otimes \cF_{\nu}
      \to \mathbb{C}
    \]
    in the Fock-space representation $\cF_\lambda$ of $G_v$, where
    $\cF_\lambda$ is the Fock module labeled by the partition $\lambda$.
  \item \textbf{The shadow}: the $\mathbb{C}^3$ partition function
    $Z(\mathbb{C}^3) = C_{\emptyset\emptyset\emptyset}(q) = M(q)
    = \prod_{n=1}^\infty (1-q^n)^{-n}$ is the denominator identity of $G_v$.
\end{enumerate}
\end{construction}

\begin{remark}[MacMahon and DT]
The MacMahon function $M(q) = \sum_{n=0}^\infty p_3(n)\, q^n$, which counts
3D partitions, is the DT partition function of $\mathbb{C}^3$ (by
the DT/GW correspondence of Maulik--Nekrasov--Okounkov--Pandharipande).
Its role as the denominator identity of the local root datum of $G(\mathbb{C}^3)$
is the simplest instance of Theorem~CY-B.
\end{remark}

\subsection{Gluing = tensor product of root data}
\label{sub:gluing}

The factorization \eqref{eq:top-vertex-Z} has a clean interpretation in
terms of root data:

\begin{construction}[Gluing of root data]
\label{constr:gluing}
The global root datum $\mathcal{R}(X)$ is assembled from local root data
$\mathcal{R}_v$ by the following procedure.

\begin{enumerate}[label=(\roman*)]
  \item \textbf{At each vertex $v$}: the local root datum $\mathcal{R}_v =
    \mathcal{R}(\mathbb{C}^3)$ with its vertex amplitude $C_{\lambda\mu\nu}$.
  \item \textbf{At each edge $e$ joining vertices $v, w$}: the ``propagator''
    is the edge weight $(-Q_e)^{|\lambda_e|} f_{\lambda_e}(q)$, which
    implements the contraction
    \[
      \cF^{(v)}_{\lambda_e} \otimes (\cF^{(w)}_{\lambda_e})^*
      \xrightarrow{\;\mathrm{contract}\;} \mathbb{C},
    \]
    i.e., the identification (tensor product pairing) of the Fock-space
    representation of $G_v$ with the dual of the Fock-space representation
    of $G_w$ along the shared edge.
  \item \textbf{The lattice}: $\Lambda(X) = \bigoplus_e \mathbb{Z} \cdot [e]$
    (one generator per internal edge = one K\"ahler class), plus the
    embedding $\Lambda_v \hookrightarrow \Lambda(X)$ at each vertex.
  \item \textbf{The global quantum group}: $G(X)$ is the tensor product
    (over the toric diagram) of the local $G_v$'s, quotiented by the
    edge identifications.  In the language of factorization algebras,
    this is the factorization product over the toric diagram.
\end{enumerate}

The summation over partitions $\lambda_e$ in \eqref{eq:top-vertex-Z}
is the ``path integral'' over the internal degrees of freedom: it sums
over all representations propagating along each edge, weighted by the
K\"ahler factor $Q_e^{|\lambda_e|}$.  This is the trace over the tensor
product:
\begin{equation}\label{eq:gluing-trace}
  Z(X) = \mathrm{Tr}_{\bigotimes_e \cF_e}
    \left( \prod_v \Phi_v \cdot \prod_e (-Q_e)^{L_e^{(0)}} q^{L_e} \right),
\end{equation}
where $L_e^{(0)}$ and $L_e$ are the zero-mode and total energy operators on the
edge Fock spaces.
\end{construction}

\begin{remark}[Relation to the crystal melting picture]
The factorization over vertices is the algebraic incarnation of the crystal
melting picture of Okounkov--Reshetikhin--Vafa: 3D partitions label states
in the Fock space, the vertex operator creates/annihilates boxes, and the
partition function counts weighted crystal configurations.  The root datum
language merely repackages this as data of the quantum vertex chiral group.
\end{remark}

\begin{conjecture}[Topological vertex as chiral intertwiner]
\label{conj:vertex-intertwiner}
The topological vertex $C_{\lambda\mu\nu}(q)$ coincides with the matrix
element of the $E_2$-chiral intertwiner of $G(\mathbb{C}^3) \simeq Y(\widehat{\mathfrak{gl}}_1) \simeq \mathcal{W}_{1+\infty}$:
\[
  C_{\lambda\mu\nu}(q) = \langle \lambda |
    \Phi^{E_2}_{A_{\mathbb{C}^3}}(z_1, z_2, z_3) | \mu, \nu \rangle,
\]
where $\Phi^{E_2}_{A_{\mathbb{C}^3}}$ is the genus-0, 3-point amplitude of
the $E_2$-chiral algebra $A_{\mathbb{C}^3} = \Phi(\cC_{\mathbb{C}^3})$,
and the Fock states $|\lambda\rangle$ are identified with basis elements
of the bar complex $B(A_{\mathbb{C}^3})$.
\end{conjecture}


%% ================================================================
\section{Connection to the Volume I shadow obstruction tower}
\label{sec:shadow-tower}

\subsection{The leading term: $F_g = \kappa \cdot \lambda_g$}
\label{sub:leading-term}

Volume~I establishes that for any chiral algebra $A$ with modular
characteristic $\kappa(A)$, the genus-$g$ obstruction on the
\emph{uniform-weight lane} (the component of highest Hodge weight in
$H^\bullet(\Mbar_g)$) is:
\[
  \mathrm{obs}_g(A)\big|_{\mathrm{uniform}} = \kappa(A) \cdot \lambda_g,
\]
where $\lambda_g$ is the top Chern class of the Hodge bundle
$\mathbb{E}_g \to \Mbar_g$.  For the quantum vertex chiral group $G(X)$
of a CY3, the modular characteristic is:
\[
  \kappa(G(X)) = \frac{\chi(X)}{2},
\]
where $\chi(X)$ is the topological Euler characteristic (which, for a CY3,
equals $2(h^{1,1} - h^{2,1})$).

This gives the leading term of the genus-$g$ topological string amplitude:
\begin{equation}\label{eq:fg-leading}
  F_g\big|_{\mathrm{leading}} = \frac{\chi(X)}{2} \cdot
    \int_{\Mbar_g} \lambda_g \cdot (\text{known intersection number}).
\end{equation}

For $g \geq 2$, the Faber--Pandharipande formula gives:
\[
  \int_{\Mbar_{g,1}} \psi_1^{2g-2}\, \lambda_g
  = \frac{2^{2g-1}-1}{2^{2g-1}} \cdot \frac{|B_{2g}|}{(2g)!},
\]
where $B_{2g}$ is the Bernoulli number.  The constant-map contribution to
$F_g$ (the $\beta = 0$ term) is indeed:
\begin{equation}\label{eq:fg-constant}
  F_g^{\beta=0} = \frac{(-1)^{g-1}\, \chi(X)}{2}
    \cdot \frac{|B_{2g}|\, |B_{2g-2}|}{2g\,(2g-2)\,(2g-2)!},
\end{equation}
which is precisely $\kappa(G(X)) \cdot \lambda_g$ evaluated via the
Mumford relation.

\subsection{Shadow corrections beyond the leading term}
\label{sub:shadow-corrections}

The full genus-$g$ amplitude $F_g = \sum_\beta N_{g,\beta}\, Q^\beta$
includes contributions from all curve classes $\beta \neq 0$.
These are the \emph{shadow corrections} to the leading term:

\begin{equation}\label{eq:fg-shadow}
  F_g = \underbrace{\kappa(G(X)) \cdot \lambda_g}_{\text{uniform-weight lane}}
    + \underbrace{\sum_{\beta \neq 0} N_{g,\beta}\, Q^\beta}_{\text{shadow corrections}}.
\end{equation}

The shadow corrections are organized by the root lattice of $G(X)$:
\begin{enumerate}[label=(\roman*)]
  \item The $\beta$-th Fourier coefficient $N_{g,\beta}$ is the genus-$g$
    amplitude of the root $\beta$ in the shadow obstruction tower.  For $g = 0$, these
    are the tree-level structure constants (Yukawa couplings).  For $g = 1$,
    they include the genus-1 one-loop correction controlled by the
    Ray--Singer torsion.
  \item The $Q^\beta$ fugacity is the root-lattice fugacity: $Q^\beta =
    e^{-\langle t, \beta \rangle}$ where $t$ is the point in the K\"ahler
    moduli space (= the Cartan element of $G(X)$).
  \item The shadow obstruction tower $\Theta_{A_X}^{\leq r}$ truncated at
    arity $r$ captures the $F_g$ up to corrections from roots of
    ``arity $> r$'' (roughly, those $\beta$ whose length in the root lattice
    exceeds $r$).
\end{enumerate}

\begin{remark}[The three towers]
Three expansions of the topological string free energy are unified:
\begin{enumerate}
  \item \textbf{Genus expansion}: $F = \sum_g F_g\, g_s^{2g-2}$.  This is
    the shadow obstruction tower graded by genus.
  \item \textbf{Instanton expansion}: $F_g = \sum_\beta N_{g,\beta}\, Q^\beta$.
    This is the root-lattice decomposition of the genus-$g$ shadow.
  \item \textbf{BPS expansion}: $F = \sum_{g,\beta,k} \frac{n^g_\beta}{k}
    (2\sin \frac{kg_s}{2})^{2g-2} Q^{k\beta}$.  This is the
    Gopakumar--Vafa reformulation, which in the root datum language gives the
    multiplicities $n^g_\beta$ as the intrinsic BPS content of each root.
\end{enumerate}
\end{remark}

\subsection{The partition function as denominator identity}
\label{sub:denominator}

For the total partition function $Z = e^F$, the product structure of the GV
formula gives:
\begin{equation}\label{eq:z-product}
  Z = M(q)^{\chi(X)/2} \cdot \prod_{\beta > 0}
    \prod_{k=1}^{\infty} \prod_{j=-\infty}^{\infty}
    \left(1 - q^{j+k/2}\, Q^{k\beta}\right)^{(-1)^{2j+1}\, n^{|j|}_\beta},
\end{equation}
(schematically), which is a Weyl--Kac--Borcherds-type infinite product ---
the denominator identity of the quantum vertex chiral group $G(X)$.

\begin{itemize}
  \item The leading factor $M(q)^{\chi(X)/2}$ is the MacMahon function
    raised to the modular characteristic $\kappa(G(X))$.  This is the
    contribution of the Weyl vector (constant maps).
  \item The product over roots $\beta > 0$ with multiplicities $n^g_\beta$
    is the automorphic correction: the imaginary roots of $G(X)$ with
    their BPS multiplicities.
  \item For $K3 \times E$: $Z$ is related to $(\Delta_5)^{-2}$, and the
    product formula is the Borcherds denominator identity of
    $\mathfrak{g}_{\Delta_5}$.
  \item For toric CY3: $Z$ is the DT partition function, factorized via the
    topological vertex.
\end{itemize}


%% ================================================================
\section{Summary: the complete dictionary}
\label{sec:dictionary}

We collect the full dictionary between the topological string on a CY3 $X$
and the quantum vertex chiral group $G(X)$.

\begin{center}
\renewcommand{\arraystretch}{1.5}
\begin{tabular}{lll}
  \hline
  \textbf{Topological string} & \textbf{$G(X)$ / Vol III} & \textbf{Vol I} \\
  \hline
  String coupling $g_s$ & Quantization parameter $\hbar$ & $\hbar$ \\
  K\"ahler parameter $Q^\beta$ & Root-lattice fugacity & --- \\
  Genus-$g$ amplitude $F_g$ & Genus-$g$ obstruction & $\mathrm{obs}_g(A_X)$ \\
  GW invariant $N_{g,\beta}$ & Shadow amplitude at $\beta$ & $\Theta^{(g)}|_\beta$ \\
  GV invariant $n^g_\beta$ & BPS root multiplicity & Intrinsic root content \\
  BCOV propagator $S^{ij}$ & Homotopy transfer kernel & Degeneration data \\
  Yukawa $C_{ijk}$ & Cubic shadow (arity 3) & $C$ in $\Theta^{\leq 3}$ \\
  Hol.\ anomaly equation & MC at genus $g$ & $D\Theta^{(g)} + \cdots = 0$ \\
  Holomorphic ambiguity & MC gauge freedom & Shadow gauge \\
  $F_g^{\beta=0}$ (constant maps) & $\kappa \cdot \lambda_g$ & Uniform-weight lane \\
  $F_g^{\beta \neq 0}$ (instantons) & Shadow corrections & Off-diagonal shadow \\
  Partition function $Z = e^F$ & Denominator identity & Bar-complex Euler product \\
  Topological vertex $C_{\lambda\mu\nu}$ & Local root datum & $E_2$-chiral intertwiner \\
  Edge gluing & Tensor product of root data & Factorization product \\
  Mirror symmetry & HMS: $\Fuk(X) \simeq D^b(\Coh(X^\vee))$ & Koszul duality \\
  DT/GW correspondence & Cohomological Hall $\leftrightarrow$ GW & $E_1 \hookrightarrow E_2$ \\
  \hline
\end{tabular}
\end{center}


%% ================================================================
\section{Open questions and future directions}
\label{sec:open}

\begin{enumerate}[label=\textbf{Q\arabic*.}]
  \item \textbf{Full proof of Conjecture~\ref{conj:gw-shadow}.}  The
    identification of GW invariants with shadow amplitudes is established at
    the level of the leading term ($\kappa \cdot \lambda_g$) and for toric
    CY3 (via the topological vertex).  The general case requires constructing
    the virtual fundamental class from the shadow obstruction tower --- this is the
    ``quantization'' step in the CY-to-chiral functor.

  \item \textbf{Refined topological vertex and motivic root data.}  The
    refined topological vertex $C_{\lambda\mu\nu}(q,t)$ (Iqbal--Kozcaz--Vafa),
    which computes the motivic DT invariants, should correspond to the
    motivic refinement of the root datum --- replacing integer multiplicities
    by classes in the Grothendieck ring of varieties.

  \item \textbf{Compact CY3 and the holomorphic anomaly.}  For compact CY3,
    the BCOV equation is essential (unlike for toric CY3, where the B-model
    is a free theory and $F_g$ is holomorphic for $g \geq 2$ in appropriate
    coordinates).  Making Construction~\ref{constr:bcov-mc} fully rigorous
    requires identifying the Costello--Li BCOV theory with the chiral bar
    dgLa of $G(X)$.

  \item \textbf{Non-perturbative completion.}  The genus expansion $F = \sum_g
    F_g\, g_s^{2g-2}$ is asymptotic (the $F_g$ grow as $(2g)!$).  The
    non-perturbative completion should come from the resurgent structure of the
    shadow obstruction tower, with Borel singularities corresponding to the real roots of
    $G(X)$ (large-radius instantons).

  \item \textbf{The $K3 \times E$ topological string.}  The topological
    string on $K3 \times E$ is not toric, so the topological vertex does not
    directly apply.  However, the BKM superalgebra $\mathfrak{g}_{\Delta_5}$
    provides the root datum, and the Igusa cusp form $\Delta_5$ is the
    denominator identity.  Constructing $F_g$ from the shadow obstruction tower in this
    case should recover the Oberdieck--Pixton formula
    (cf.\ the $K3 \times E$ chapter of Volume~III).

  \item \textbf{Relation to wall-crossing.}  The DT/GW correspondence and the
    wall-crossing formula of Kontsevich--Soibelman should be visible in the
    root datum as changes of Weyl chamber --- the shadow obstruction tower $\Theta_{A_X}$
    should be wall-crossing invariant, with its chamber-dependent expansion
    recovering the different BPS spectra.
\end{enumerate}

%% ================================================================
\section{The MacMahon obstruction}
\label{sec:macmahon-obstruction}
%% ================================================================

The MacMahon function $M(q) = \prod_{n \geq 1} (1-q^n)^{-n}$
is the DT partition function of $\mathbb{C}^3$ and the denominator
identity of the $\mathbb{C}^3$ quantum vertex chiral group.
A natural question: does the shadow obstruction tower of
$A_{\mathbb{C}^3}$ reproduce $M(q)$ term by term?

The swarm computation (\texttt{compute/lib/macmahon\_shadow\_decomposition.py},
50~tests) gives a precise negative answer.

\subsection{The mismatch}

Write the genus expansion of both objects.  The topological string
free energy of $\mathbb{C}^3$ is $F = \sum_{g \geq 0} F_g \, g_s^{2g-2}$,
and the shadow obstruction tower has projections
$\Theta^{(g)} = \kappa \cdot \lambda_g + \text{(higher arity)}$.
At each genus $g$, the topological string coefficient is
\[
  F_g^{\mathrm{top}}(\mathbb{C}^3) \;=\;
  \frac{(-1)^{g-1} B_{2g}}{2g(2g-2)!}
  \qquad (g \geq 2),
\]
involving a \emph{single} Bernoulli number $B_{2g}$.  The
shadow obstruction tower, by contrast, produces
\[
  F_g^{\mathrm{sh}} \;=\; \kappa \cdot \lambda_g^{FP}
  \;=\; \kappa \cdot
  \frac{(-1)^{g-1} B_{2(g-1)} \cdot B_{2g}}{(2g)! \cdot (2g-2)}
  \qquad (g \geq 2),
\]
involving the \emph{product} $B_{2(g-1)} \cdot B_{2g}$
(the Faber--Pandharipande class).  These differ: the shadow
tower has a double Bernoulli product where the topological
string has a single Bernoulli number.

\subsection{Generating-function-level agreement}

Despite the genus-by-genus mismatch, the two objects agree at
the level of generating functions.  Let
\[
  \widehat{A}(x) = \frac{x/2}{\sinh(x/2)}
  = 1 - \frac{x^2}{24} + \frac{7x^4}{5760} - \cdots
\]
be the $\hat{A}$-genus.  Then:
\begin{itemize}
  \item The shadow tower generating function is
    $\sum_{g \geq 1} F_g^{\mathrm{sh}} \, \hbar^{2g}
    = \kappa \cdot (\widehat{A}(i\hbar) - 1)$.
  \item The topological string generating function
    $\sum_{g \geq 2} F_g^{\mathrm{top}} \, g_s^{2g-2}$
    is the asymptotic expansion of $\log M(e^{-g_s})$.
\end{itemize}
The relation between the two is through the Laplace/Borel
transform: the double Bernoulli product $B_{2(g-1)} B_{2g}$
is the convolution of two single-Bernoulli series.  The
$\hat{A}$-genus packages this convolution into a closed form
that is \emph{not} the same as the MacMahon expansion genus
by genus, but whose Borel sum agrees.

\subsection{The obstruction}

\begin{warning}[MacMahon $\neq$ scalar shadow, term by term]
$M(q)$ is NOT the scalar shadow tower evaluated term by term.
The scalar shadow $F_g^{\mathrm{sh}} = \kappa \lambda_g^{FP}$
differs from $F_g^{\mathrm{top}}(\mathbb{C}^3)$ at every
genus $g \geq 2$.  The agreement is at the level of generating
functions (Borel-summed), not at the level of individual
Taylor coefficients.
\end{warning}

\noindent
The physical interpretation: the shadow tower computes the
\emph{algebraic} genus expansion (from the bar complex on
$\overline{\mathcal{M}}_{g,n}$), while the MacMahon function
comes from the \emph{analytic} partition function (the
regularized Euler product on the root lattice).  The passage
from algebraic to analytic requires a resummation that
reshuffles genus-by-genus coefficients.

\bibliographystyle{alpha}
% References are inline for this note.

\end{document}
